% Section 10.3, from lectures/L10/README.md: the objectives, the evaluation questions, and what comes
% after the last lecture.

\needspace{10\baselineskip}
\section{Review}
\label{c10:sec:review}

\subsection*{What you should be able to do now}
After this chapter, you should:
\begin{itemize}
  \item Have created a concrete subclass \code{Flatten} that inherits
    \code{ml::flatten_layer::Interface}.
  \item Have implemented \code{feedforward()} and \code{backpropagate()} for \code{Flatten} (note:
    no \code{optimize()}; the interface has none, since a flatten layer has no trainable
    parameters).
  \item Have a fully working, end-to-end CNN:
    \code{input → Conv → MaxPool → Flatten → Dense → prediction}.
  \item Recognize the shared feedforward, backpropagate and optimize pattern common to every layer
    type built in this book.
  \item Be able to compare the layer types built in this book (\code{Fixed}, \code{Dense},
    \code{Conv}, \code{MaxPool}, \code{Flatten}) in terms of trainable parameters, computation, and
    role in the network.
  \item Be able to reason about memory footprint, numeric precision, and real-time constraints when
    deploying a trained network outside of a general-purpose computer.
\end{itemize}

\needspace{6\baselineskip}
\subsection*{Questions to test yourself}
You should be able to explain:
\begin{enumerate}
  \item Can you trace a single training example through your finished \code{cnn_work} pipeline,
    from feedforward through backpropagation to weight updates, describing what happens at each
    layer?
  \item What do all the layer types built in this book have in common in their interface
    (feedforward, backpropagate, optimize), and where do they genuinely differ?
  \item Which layer types have trainable parameters, and which don't? Why?
  \item If you were designing a new CNN for a different image classification task, how would you
    decide how many conv and pooling layers to use, and why?
  \item Why might a network that runs comfortably on a laptop fail to fit, or fail to meet a
    deadline, on a microcontroller?
  \item What's the trade-off between fixed-point and floating-point arithmetic, and why does it
    matter more on some targets than others?
\end{enumerate}

\needspace{6\baselineskip}
\subsection*{Looking ahead}
This is the last chapter, and with it the book's main text is complete. What remains is to work the
exercises in \secref{c10:sec:exercises}, and to compare the Calculation and Reflection exercises
with their worked solutions in Appendix~\ref{sol:ch:solutions}. After that, the two written papers
in Appendices~\ref{exam:ch:about} to \ref{ansb:ch:answers} are there to test, on paper and without a
compiler, how much of all ten chapters you can reconstruct on your own.
